\section{The Join You Cannot See}
\label{sec:ch05-the-join-you-cannot-see}

\index{join!inside one process|(}%
Start the service and ask for the schedule with its speakers.

\begin{minted}{text}
POST /graphql HTTP/1.1
Host: localhost:5001
Content-Type: application/json

{"query":"{ sessions(first: 10) { nodes { title speaker { name } } } }"}
\end{minted}

\begin{minted}[fontsize=\footnotesize]{json}
{"data":{"sessions":{"nodes":[
  {"title":"Schemas That Outlive Their Authors","speaker":{"name":"Ada Fischer"}},
  {"title":"Reading a Query Plan","speaker":{"name":"Ada Fischer"}},
  {"title":"Paging Without Offsets","speaker":{"name":"Bruno Kaminski"}},
  {"title":"The Bank That Split Its Graph","speaker":{"name":"Chidi Okafor"}}]}}}
\end{minted}

\index{statement count!the seam query}%
Two statements, as chapter~\ref{ch:data-without-n-plus-1} left it. The second
one is the one to look at.

\begin{minted}{text}
-- statement 2
SELECT "s"."Id", "s"."Bio", "s"."Name"
FROM "Speakers" AS "s"
WHERE "s"."Id" IN (@ids1, @ids2, @ids3)
\end{minted}

\index{key!carried by the foreign key}%
Three parameters for four sessions, because Ada Fischer gives two of them.
Those three values are the join. They came off the \code{SpeakerId} column of
the four rows the first statement returned, went into a DataLoader, and came
back out as a set. Every one of them is an integer that lived for the length of
one request inside one process.

\index{response!carries no key}%
Now read the response again and look for them. They are not there. Neither is
anything else that ties a speaker to a session: no \code{speakerId}, no
\code{sessionId}, no key of any kind. The nesting is the whole record of the
join, and a client that flattened this document into two lists would destroy
information it could not recover.

\index{SpeakerId@\code{SpeakerId}!absent from the schema by design}%
That is not an accident of what this query asked for. There is no field
anywhere on \code{Session} that would return it.
Chapter~\ref{ch:data-without-n-plus-1} put \code{[GraphQLIgnore]} on
\code{Session.SpeakerId} on purpose, and chapter~\ref{ch:schema-design} took
the other integer key out of the schema as well. A reader who typed this
service out of the book has, deliberately, no way to ask it which speaker
belongs to which session except by asking for the speaker.

\index{join!is a private implementation detail}%
So the join is real, it is cheap, and it is invisible from outside. It is also,
right now, a private matter between \code{SessionType.cs} and a table. Nothing
in the schema commits you to performing it, nothing in the schema describes
how, and you could swap the DataLoader for a dictionary this afternoon without
telling anyone.

\index{seam!where the line goes}%
Draw the line. Speakers belong to another team now: another repository, another
deployment, another database that your process cannot open a connection to.
The first statement is still yours. The second one is not, and the resolver
that would have issued it is holding three integers with nothing to do.

\index{join!has to move outward}%
There are only three places that join can go. It can go into every client,
which is the answer chapter~\ref{ch:do-you-need-this} spent a chapter arguing
against, and which the schema you have just built cannot even support. It can
stay in your service, which means your service calls theirs, which is a single
GraphQL service with extra steps and the same coordination problem underneath
it. Or it can move outward, into something in front of both services that
belongs to neither.

Figure~\ref{fig:ch05-the-join-moves} is the third answer, drawn beside the
arrangement you have now.

\begin{figure}
  \centering
  % The same join, performed twice: once inside one process, once across a seam.
%
% Same idiom as figures/tikz/ch01-one-field.tex, ch02-two-requests.tex and
% ch03-one-batch.tex: each lane is a timeline read left to right, ticks mark
% stages, a dashed vertical marks a boundary the work crosses, and a brace
% under the stages that matter carries a one-line label. Lane A is the
% arrangement chapter 3 finished with; lane B is the same join with a router
% in the middle. The four ticks sit at the same x in both lanes, because the
% argument of the figure is that the two arrangements do the same work in the
% same order and differ only in what the keys have to travel through, so any
% horizontal offset between the lanes would be saying something false.
%
% This figure fixes two visuals SPEC decision 33 says are reused. The seam is
% a dashed vertical rule; one round trip out and back crosses it twice, so
% lane B carries two rules and not four. The router is drawn as the actor
% whose lane it is, named in the tick labels, rather than as a box, because a
% box would make it look like a fourth service rather than the thing holding
% the plan.
%
% Bare tikzpicture (house style): the figure environment, caption and label
% live at the call site in the section file.
\begin{tikzpicture}[
  x=1mm, y=1mm,
  every node/.append style={font=\sffamily\scriptsize},
  axis/.style={draw, line width=0.4pt, -{Straight Barb[length=1.4mm, width=1.6mm]}},
  tick/.style={draw, line width=0.4pt},
  event/.style={anchor=south, align=center, text width=24mm, inner sep=1pt},
  span/.style={draw, line width=0.4pt},
  spanlabel/.style={anchor=north, align=center, inner sep=2pt},
  lane/.style={anchor=west, font=\sffamily\scriptsize\itshape},
  boundary/.style={draw, line width=0.4pt, dashed},
]

% ------------------------------------------------------------------- lane A
\node[lane] at (0,14) {A. one service, one process: the join you have now};

\draw[axis] (0,0) -- (150,0);
\foreach \x in {16,58,100,138}{\draw[tick] (\x,-1.5) -- (\x,1.5);}

\node[event, text width=26mm] at (16,3)  {sessions resolver:\\statement 1, 4 sessions};
\node[event, text width=26mm] at (58,3)  {speaker resolver x4:\\3 distinct keys held};
\node[event, text width=26mm] at (100,3) {statement 2: the 3\\keys in one IN list};
\node[event, text width=20mm] at (138,3) {one response\\assembled};

\draw[span] (50,-4) -- (50,-7) -- (108,-7) -- (108,-4);
\node[spanlabel] at (79,-8) {three integers, held in memory,\\named in no schema anywhere};

% ------------------------------------------------------------------- lane B
\begin{scope}[shift={(0,-50)}]
  \node[lane] at (0,14) {B. two subgraphs and a router: the same join, the same order};

  \draw[axis] (0,0) -- (150,0);
  \foreach \x in {16,58,100,138}{\draw[tick] (\x,-1.5) -- (\x,1.5);}
  \foreach \x in {37,119}{
    \draw[boundary] (\x,-2) -- (\x,8);
    \node[spanlabel, inner sep=1pt] at (\x,-2) {seam};
  }

  \node[event, text width=26mm] at (16,3)  {router asks sessions\\for the page};
  \node[event, text width=26mm] at (58,3)  {sessions answers: a\\speaker key per session};
  \node[event, text width=26mm] at (100,3) {router asks speakers,\\the 3 keys in one call};
  \node[event, text width=20mm] at (138,3) {router merges,\\one response};

  \draw[span] (50,-6) -- (50,-9) -- (108,-9) -- (108,-6);
  \node[spanlabel] at (79,-10) {the same three keys, serialized, through\\a field both services have to declare};
\end{scope}

\end{tikzpicture}

  \caption{The same join, twice. Lane A is the service as
  chapter~\ref{ch:data-without-n-plus-1} left it: the three keys are collected
  and spent without ever being named in a schema. Lane B does the same work in
  the same order, with two round trips through the middle of it. Nothing about
  the join changed, and that is the argument: what changed is that the keys now
  have to be serialized, which means they have to be fields, which means both
  services have to agree on what they are called and what is in them.}
  \label{fig:ch05-the-join-moves}
\end{figure}

\index{key!has to become public}%
Read the two lanes for what is different, which is not the number of steps. In
lane A three integers move between two lines of C\# in the same call stack. In
lane B the same three values have to survive being written to JSON, sent over a
network, parsed by a service that has never seen your database, and matched
against rows in a table you do not own. A value that has to do all of that is
not an implementation detail any more. It is an API.

\index{key!is the whole of the model}%
So the rest of this chapter is about who is allowed to be on each side of that
boundary, and what each of them has to declare before the value may cross.
\index{join!inside one process|)}%
